\subsubsection{Bijlagen}
\label{subsubsec:bijlagen}

Een bijlage kan schadelijke code uitvoeren op het toestel van het slachtoffer, zodra het bestand wordt geopend. 
Detectiekenmerken voor bijlagen zijn te vinden in de bestandsnaam en metadata, alsook in de statische documentkenmerken. Deze kenmerken leveren informatie op zonder dat de bijlage wordt uitgevoerd.

\begin{enumerate}
    \item \textbf{Dubbele bestandsextensie.} Aanvallers kunnen een uitvoerbaar bestand verbergen achter een onschuldig ogende extensie, zoals \emph{example.jpg.exe} \autocite{Alazab2016}. Enkel de laatste extensie bepaalt het bestandstype.
    \item \textbf{Metadata-kenmerken.} Het MIME-type geeft aan om welk soort bestand het gaat. Een discrepantie tussen het MIME-type en de extensie kan erop wijzen dat een bijlage een ander bestandstype is dan gedacht \autocite{Mendes2025}. Ook de bestandsgrootte kan als detectiekenmerk dienen, bijvoorbeeld bij bijlagen met een ongewoon grote omvang \autocite{Mendes2025}.
    \item \textbf{Statische documentkenmerken.} Bij PDF- en Office-documenten kan de interne structuur worden onderzocht op kwaadaardige JavaScript-code \autocite{Srndic2013}. Hiervoor hoeft het document niet geopend te worden.
\end{enumerate}

Kenmerken van bijlagen kunnen bruikbare signalen zijn voor phishingdetectie, maar blijken uit de experimenten van \textcite{Mendes2025} geen doorslaggevende factor te zijn. De toevoeging van bijlagekenmerken aan een model gebaseerd op tekstkenmerken, heeft geleid tot een verbetering van slechts 0.2 procentpunt.
